Four dimensions of democratic representation quality have been established in
this track: electoral competitiveness (O.1), voter turnout (O.2), legislative
polarization (O.3), and constituent-representative geographic distance (O.4).
Each dimension independently favours algorithmic redistricting over partisan
gerrymandering.
This paper constructs a \emph{Representation Quality Index} (RQI) that
synthesises all four dimensions into a single 0--1 composite score.

The RQI assigns equal weight to each dimension (sensitivity analysis shows
rankings are robust to weight variation, as with K.7's compactness composite).
Under equal weights, \textsc{bisect} algorithmic plans score $0.62 \pm 0.04$
on the 0--1 RQI scale, compared to $0.41 \pm 0.07$ for enacted partisan
gerrymanders and $0.58 \pm 0.05$ for bipartisan incumbent-protection plans
(Iowa-style neutral redistricting).

The RQI provides a court-admissible single-metric comparison for redistricting
plans that goes beyond compactness (K.7) and partisan fairness (L.1--L.6)
to encompass actual democratic outcomes.
A plan at RQI $< 0.50$ is below the midpoint of the democratic-quality scale;
all five study state enacted plans score below 0.50.
All five \textsc{bisect} plans score above 0.55.

The RQI is validated against Pew Research state-level public satisfaction
with redistricting (r = 0.54 with RQI across 18 states with available data),
confirming that the composite captures public perceptions of redistricting
quality.
